\documentclass[../Main.tex]{subfiles}

\begin{document}
\section{Introduction}
\subsection{Minimisation Problems}
\begin{definition}{Optimisation problem}
    An \underline{optimisation problem} has the form:
    \begin{align*}
        \text{Minimise } & f(\vec{x}) \\
        \text{Subject to } & \vec{h}(\vec{x}) = \vec{b}, \\
        & \vec{x} \in \Xset
    \end{align*}
    where $f : \R^n \mapsto \R$ is the \underline{objective function}, $\vec{h}(\vec{x}) = \vec{b}$ is a \underline{functional constraint} where $\vec{h} : \R^n \mapsto \R^m$ and $\vec{b} \in \R^m$, and $\vec{x} \in \Xset$ is a \underline{regional constraint}.
\end{definition}
\begin{definition}{Feasible set}
    A \underline{feasible set} is the set of $\vec{x}$ that satisfy all constraints in the minimisation problem. We denote it:
    \begin{equation*}
        \Xset(\vec{b}) = \subsetselect{\vec{x} \in \Xset}{\vec{h}(\vec{x}) = \vec{b}}
    \end{equation*}
\end{definition}
\begin{definition}{Optimal value}
    The \underline{optimal value} for a minimisation problem is the vector $\vec{x}^*$ that satisfies the constraints and minimises $f$.
\end{definition}
\subsection{Defining Convexity}
\begin{definition}{Convex set}
    A set $S \subseteq \R^n$ is \underline{convex} if, for all $\vec{x}, \vec{y} \in S$ and $\lambda \in [0,1]$,
    \begin{equation*}
        \lambda \vec{x} + (1-\lambda) \vec{y} \in S
    \end{equation*}
\end{definition}
\begin{definition}{Convex function}
    A function $f : S \mapsto \R$ is \underline{convex} if $S$ is a convex set and, for all $\vec{x}, \vec{y} \in S$, $\lambda \in [0, 1]$,
    \begin{equation}
        f(\lambda \vec{x} + (1-\lambda) \vec{y}) \leq \lambda f(\vec{x}) + (1-\lambda) f(\vec{y})
        \label{eqnConvexDef}
    \end{equation}
\end{definition}
\begin{theorem}[First-order condition for convexity]
    A differentiable function $f : S \mapsto \R^n$ is convex if and only if:
    \begin{equation}
        f(\vec{y}) \geq f(\vec{x}) + \left(\nabla f(\vec{x})\right) \cdot (\vec{y} - \vec{x})
        \label{eqnConvexFirstOrder}
    \end{equation}
    for all $\vec{x}, \vec{y} \in S$.
    \label{thmConvexFirstOrder}
\end{theorem}
\begin{proof}
    Fix $\vec{x}$ and $\vec{y}$ and consider the restricted function of $t \in [0, 1]$:
    \begin{equation*}
        f((1-t)\vec{x} + t\vec{y})
    \end{equation*}
    \begin{proofdirection}{$\Rightarrow$}{Suppose $f$ is convex}
        Take the convexity definition in equation~\ref{eqnConvexDef} and rearrange as an inequality for $f(\vec{y})$. Take the limit $t \to 0$ for equation~\ref{eqnConvexFirstOrder}.
    \end{proofdirection}
    \begin{proofdirection}{$\Leftarrow$}{Suppose $f$ satisfies the first-order condition}
        For ease define $\vec{x_t} = (1-t)\vec{x} + t\vec{y}$. Apply \eqnref{eqnConvexFirstOrder} twice to get:
        \begin{align*}
            f(\vec{y}) &\geq f(\vec{x_t}) + \left[\nabla f(\vec{x_t})\right]^T (\vec{y} - \vec{x_t}) \\
            f(\vec{x}) &\geq f(\vec{x_t}) + \left[\nabla f(\vec{x_t})\right]^T (\vec{x} - \vec{x_t})
        \end{align*}
        Then a convex combination of these equations gives \eqnref{eqnConvexDef}.
    \end{proofdirection}
\end{proof}
\begin{corollary}
    Suppose $\vec{x}$ is a point where $\nabla f = \vec0$. Then $\vec{x}$ is a global minimum of the function.
    \label{corLocalMinIsGlobal}
\end{corollary}
\begin{definition}{Hessian matrix}
    The \underline{Hessian matrix} of a function $f : \R^n \mapsto \R$ is the symmetric matrix $H(f)$ with elements:
    \begin{equation*}
        \left[H(f)\right]_{ij} = \frac{\partial^{2}f}{\partial x_i\partial x_j}
    \end{equation*}
\end{definition}
\begin{definition}{Positive semi-definiteness}
    A matrix $A \in R^{n \times n}$ is \underline{positive semi-definite} if, for all $\vec{v} \in \R^n$,
    \begin{equation*}
        \vec{x}^T A \vec{x} \geq 0.
    \end{equation*}
    In other courses we prove that this is equivalent to having all eigenvalues non-negative. We write $A \succeq 0$.
\end{definition}
With this notation we can write $A \succeq B$ if $A-B$ is positive semi-definite.
\begin{theorem}[Second-order condition for convexity]
    A twice-differentiable function $f : \R^n \mapsto \R$ is convex if and only if:
    \begin{equation}
        H(f) \succeq 0
        \label{eqnConvexSecondOrder}
    \end{equation}
    at every point in $\R^n$.
    \label{thmConvexSecondOrder}
\end{theorem}
\begin{proof}
    Only the forward direction is examinable, so assume \eqnref{eqnConvexSecondOrder}. Then the multivariate Taylor series (with Lagrange remainder) implies \eqnref{eqnConvexFirstOrder}.
\end{proof}
\section{Gradient Descent}
\subsection{Algorithm}
The algorithm is as follows:
\begin{enumerate}
    \item Start at an initial point $\vec{x_0}$.
    \item Choose a step size $\epsilon > 0$.
    \item Given point $\vec{x_k}$, $\vec{x_{k+1}} = \vec{x_k} - \epsilon\nabla f(\vec{x_k})$.
    \item Repeat until reaching either a predetermined maximum number of steps or reaching a small enough $\nabla f$.
\end{enumerate}
This works because at each step,
\begin{equation*}
    f(\vec{x_{k+1}}) \approx f(\vec{x_k}) - \epsilon \norm{\nabla f(\vec{x_k})}
\end{equation*}
for sufficiently small $\epsilon$.
\subsection{Choosing the Parameters}
\begin{definition}{$\beta$-smoothness}
    A twice-differentiable function $f : \R^n \mapsto \R$ is \underline{$\beta$-smooth} if, for $\beta > 0$,
    \begin{equation*}
        H(f) \preceq \beta I
    \end{equation*}
    for all $\vec{x} \in \R^n$.
\end{definition}
\begin{definition}{$\alpha$-strong convexity}
    A twice-differentiable function $f : \R^n \mapsto \R$ is \underline{$\alpha$-strongly convex} if, for $\alpha > 0$,
    \begin{equation*}
        H(f) \succeq \alpha I
    \end{equation*}
    for all $\vec{x} \in \R^n$.
\end{definition}
\begin{theorem}
    Let $f : \R^n \mapsto \R$ be twice-differentiable. Let $\vec{x}, \vec{y} \in \R^n$.

    If $f$ is $\beta$-smooth then:
    \begin{equation*}
        f(\vec{y}) \leq f(\vec{x}) + (\nabla f(\vec{x})) \cdot (\vec{y} - \vec{x}) + \frac{\beta}{2}\norm{\vec{x} - \vec{y}}^2
    \end{equation*}
    If $f$ is $\alpha$-strongly convex then:
    \begin{equation*}
        f(\vec{y}) \leq f(\vec{x}) + (\nabla f(\vec{x})) \cdot (\vec{y} - \vec{x}) - \frac{\alpha}{2}\norm{\vec{x} - \vec{y}}^2
    \end{equation*}
    \label{thmConvexityBounds}
\end{theorem}
\begin{proof}
    Consider the Taylor series with Lagrange remainder:
    \begin{equation*}
        f(\vec{y}) = f(\vec{x}) + (\nabla f(\vec{x})) \cdot (\vec{y} - \vec{x}) + \frac12 (\vec{y} - \vec{x})^T H(f)(\vec{z}) (\vec{y} - \vec{x})
    \end{equation*}
    for some $\vec{z}$ along the line between $\vec{x}$ and $\vec{y}$.

    Recall also that, for $A \preceq B$,
    \begin{equation*}
        \vec{x}^T A \vec{x} \leq \vec{x}^T B \vec{x}
    \end{equation*}
    for all $\vec{x} \in \R^n$.

    Then use the inequalities afforded by $\alpha$-strong convexity and $\beta$-smoothness for the result.
\end{proof}
\begin{corollary}
    For a function $f$ as above:
    \begin{equation*}
        f\left(\vec{x} + \beta^{-1} \nabla f(\vec{x})\right) \leq f(\vec{x}) - \frac{1}{2\beta} \norm{\nabla f(\vec{x})}^2
    \end{equation*}
    \label{corStepSizeBound}
\end{corollary}
\begin{remark}
    This gives us a step size, $\epsilon = \beta^{-1}$.
\end{remark}
\begin{corollary}
    Let $f^*$ be the minimum value of $f(\vec{x})$. Then for any $\vec{x} \in \R^n$,
    \begin{equation*}
        f^* \geq f(\vec{x}) - \frac{\norm{\nabla f(\vec{x})}^2}{2\alpha}
    \end{equation*}
    \label{corStopCriterionBound}
\end{corollary}
\begin{remark}
    This gives us the stopping condition for a required precision $\delta > 0$,
    \begin{equation*}
        \norm{\nabla f(\vec{x})}^2 < 2\alpha\delta \implies \abs{f(\vec{x}) - f^*} < \delta.
    \end{equation*}
\end{remark}
\subsection{Convergence}
\begin{theorem}
    Let $f : \R^n \mapsto \R$ be $\alpha$-strongly convex and $\beta$-smooth. The gradient descent algorithm with step size $\beta^{-1}$ satisfies, at iteration $T$,
    \begin{align*}
        f(\vec{x_T}) - f^* &\leq \left(1 - \frac{\alpha}{\beta}\right)^T (f(\vec{x_0}) - f^*) \\
        &\leq e^{-\frac{\alpha}{\beta}T} (f(\vec{x_0}) - f^*)
    \end{align*}
    where $f^*$ is the value of $f$ at the optimal point.
    \label{thmGradDescConvergence}
\end{theorem}
\begin{proof}
    By \thmref{corStepSizeBound} we have (after subtracting $f^*$):
    \begin{equation*}
        f(\vec{x_{t+1}}) - f^* \leq f(\vec{x_T}) - f^* - \frac1{2\beta} \norm{\nabla f(vec{x_T})}^2
    \end{equation*}
    Now do the same thing with \thmref{corStopCriterionBound} and combine with the above inequality to get:
    \begin{equation*}
        \left(1 - \frac{\alpha}{\beta}\right) (f(\vec{x_T}) - f^*)
    \end{equation*}
    and iterate this back for the result. For the second inequality, apply the well-known approximation for $e$.
\end{proof}
\section{Second-Order Methods}
\subsection{Newton's Method}
Newton's method is based on:
\begin{equation}
    \vec{x_{t+1}} = \vec{x_t} - \left(H(f)(\vec{x_{t}})\right)^{-1} \nabla f(\vec{x_{t}})
    \label{eqnNewtonMethod}
\end{equation}
We can remember this as Newton-Raphson on the first derivative (in the 1D case).
\begin{definition}{Matrix norm}
    The \underline{norm} of a matrix $A \in R^{n \times n}$ is defined as the smallest positive real value $a$ such that:
    \begin{equation*}
        \norm{A\vec{z}} \leq a\norm{\vec{z}}
    \end{equation*}
    for all $\vec{z} \in \R^n$. We then say $\norm{A} = a$.
\end{definition}
\begin{remark}
    For positive semi-definite matrices, this is the largest eigenvalue.
\end{remark}
\begin{definition}{$l$-Lipschitz}
    A twice-differentiable function $f : \R^n \mapsto \R$ is \underline{$l$-Lipschitz} if, for all $\vec{x}, \vec{y} \in R^n$,
    \begin{equation*}
        \norm{H(f)(\vec{x}) - H(f)(\vec{y})} \leq l\norm{\vec{x} - \vec{y}}
    \end{equation*}
\end{definition}
\begin{theorem}
    Let $f : \R^n \mapsto \R$ and suppose that it is $l$-Lipschitz and $\alpha$-strongly convex. Then Newton's Method satisfies:
    \begin{equation*}
        f(\vec{x_{t}}) - f^* \leq \frac{2\alpha^3}{l^2} \left(\frac{l}{2\alpha^2} \norm{f(\vec{x_{0}})}\right)^{2^{t+1}}
    \end{equation*}
    \label{thmNewtonConvergence}
\end{theorem}
\begin{proof}
    The proof of \thmref{thmNewtonConvergence} is non-examinable.
\end{proof}
\begin{remarks}
    \item This is exponentially faster than gradient descent
    \item This only gives us information about the convergence if $\norm{f(\vec{x_{0}})} < \frac{2\alpha^2}{l}$. Therefore, we often run gradient descent to get to this value and then run Newton's method which is faster.
\end{remarks}
\subsection{Barrier Method}
Suppose we now impose the regional constraint:
\begin{equation*}
    A \vec{x} - \vec{b} \leq 0
\end{equation*}
for $A \in \R^{m \times n}$, $\vec{b} \in \R^m$.

Then we define:
\begin{equation*}
    \Phi(\vec{x}) = -\sum_{i=1}^{m} \log(-(\vec{a_{i}}^T \vec{x} - b_i))
\end{equation*}
which is small inside the feasible region and large outside it. Our objective function then becomes $g_t(\vec{x}) = tf(\vec{x}) + \Phi(\vec{x})$, and the new Newton's Method algorithm is:
\begin{enumerate}
    \item Find a strictly feasible $\vec{x}$ (by running gradient descent on $\Phi$, for example).
    \item Set $t > 0$ and $s > 1$.
    \item Compute $\vec{x}^*(t)$ by minimising $g_t$ with Newton's method, starting with $\vec{x}$ \label{stepBarrierMethod}
    \item Update $\vec{x} = \vec{x}^*(t)$, $t = st$.
    \item If $t$ is not large enough, return to step \ref{stepBarrierMethod}.
\end{enumerate}
\end{document}